\documentclass[tikz,border=4pt]{standalone}
\usepackage{tikz}
\usetikzlibrary{arrows.meta}

% Sample 1: n = 4, k = 2. Turning on (0,0) and (2,2) lights 4 lighthouses:
% (0,0) spreads to (0,1) (same x), which spreads to (1,1) (same y).
\begin{document}
% The HTML statement renders this PDF with a transparent background and Kattis
% does not invert illustrations in dark mode, so every colour is a mid-tone that
% reads on both a white and a dark (#333) background.
\definecolor{ink}{HTML}{8C8C8C}
\definecolor{beam}{HTML}{E08A00}
\definecolor{lamp}{HTML}{FFB81C}
\definecolor{offlamp}{HTML}{7A7A7A}
\begin{tikzpicture}[
    lit/.style={circle, draw=beam, fill=lamp, minimum size=4mm, inner sep=0pt},
    off/.style={circle, draw=ink, fill=offlamp, minimum size=4mm, inner sep=0pt},
    halo/.style={circle, fill=lamp, fill opacity=0.3, minimum size=8.5mm, inner sep=0pt},
    spread/.style={-{Stealth[length=2mm]}, beam, thick, shorten <=2.4mm, shorten >=2.4mm},
  ]
  \draw[ink, opacity=0.6] (0,0) grid (4,4);
  \foreach \i in {0,...,3} {
    \node[below, ink] at (\i+0.5,0) {\small $\i$};
    \node[left, ink]  at (0,\i+0.5) {\small $\i$};
  }
  \node[ink] at (4.35,-0.3) {$x$};
  \node[ink] at (-0.3,4.35) {$y$};

  \foreach \p in {(0.5,0.5),(0.5,1.5),(1.5,1.5),(2.5,2.5)} \node[halo] at \p {};
  \draw[spread] (0.5,0.5) -- (0.5,1.5);
  \draw[spread] (0.5,1.5) -- (1.5,1.5);
  % rays mark the two lighthouses that are turned on
  \foreach \c in {(0.5,0.5),(2.5,2.5)}
    \foreach \a in {0,45,...,315}
      \draw[beam, semithick, shift={\c}] (\a:2.7mm) -- (\a:3.7mm);
  \foreach \p in {(0.5,0.5),(0.5,1.5),(1.5,1.5),(2.5,2.5)} \node[lit] at \p {};
  \node[off] at (3.5,3.5) {};
\end{tikzpicture}
\end{document}
